% Part: proof-theory
% Chapter: normalization
% Section: normalization-thm

\documentclass[../../../include/open-logic-section]{subfiles}

\begin{document}

\olfileid{pt}{nor}{thm}

\olsection{نارمل روپ وچ لیاؤن دا قضیہ}

!!^a{proof} \emph{نارمل} روپ وچ اے جے اوہدے وچ کوئی قطع ٹوٹا نہیں۔ صاف اے کہ ایہ اُدوں تے صرف اُدوں اے جدوں اسیں اوہدے اُتے نہ تھاں بدلّن والی تبدیلی لا سکدے آں تے نہ تقلیل والی۔ !!a{proof} نوں \emph{نارمل روپ وچ لیاؤن} دا مطلب اے اوہنوں تھاں بدلّن تے تقلیل والیاں تبدیلیاں لا کے نارمل ثبوت وچ بدلنا، جدوں تک کوئی قطع نہیں رہندا۔ ایہ ہمیشہ ہو سکدا اے، ایہو \emph{نارمل روپ وچ لیاؤن دا قضیہ} اے۔

\begin{thm}\ollabel{thm:normalization}
کوئی وی \Log{N1i}-!!{proof} $\delta$ نارمل روپ وچ آ سکدا اے، یعنی قطعاں توں بناں !!a{proof}~$\delta'$ وچ بدلیا جا سکدا اے۔
\end{thm}

\begin{proof}
قطع درجے تے قطع لمبائی اُتے استقرا نال،~$\delta$ دے۔ استقرا دا مفروضہ ایہ اے کہ ہر اوہ !!{proof} جیہدا قطع درجہ گھٹ اے، یا قطع درجہ اوہی پر قطع لمبائی گھٹ اے، نارمل روپ وچ آ سکدا اے۔

جے قطع لمبائی~$0$ اے، تاں کوئی قطع نہیں تے $\delta$ پہلاں ای نارمل اے۔ سو منّ لو~$\delta$ دا قطع درجہ تے قطع لمبائی~$>
0$ نیں۔ سبھ توں اُتلا تے سبھ توں سجّا زیادہ توں زیادہ درجے والا قطع ٹوٹا چُنو۔ جے قطع ٹوٹے دی لمبائی $>1$ اے، تاں تھاں بدلّن والی تبدیلی لاؤ۔ جے قطع ٹوٹے دی لمبائی~$1$ اے، تاں میل کھاندی تقلیل والی تبدیلی لاؤ۔ ایس نال اک نواں !!{proof} ملدا اے جیہدا یا قطع درجہ اوہی اے پر قطع لمبائی گھٹ اے، یا قطع درجہ گھٹ اے۔ استقرا دا مفروضہ لگدا اے۔
\end{proof}

اُتلے ثبوت دا طریقہ (ہمیشہ سبھ توں اُتلا تے سبھ توں سجّا زیادہ توں زیادہ درجے والا قطع لو) ایہ اشارہ دیندا اے کہ !!a{proof} اُتے تبدیلیاں کسے خاص ترتیب نال لاؤنیاں پین گئیاں تاں جو نارمل !!{proof} ملے۔ حیرانی دی گل ایہ اے کہ ترتیب نال کوئی فرق نہیں پیندا۔ تھاں بدلّن تے تقلیل والیاں تبدیلیاں کسے وی ترتیب نال لاؤ، آخر نارمل ثبوت مل جاندا اے۔

ٹوٹے تے قطع دی تعریف \Log{N2i} لئی وی سوکھے ای دِتّی جا سکدی اے، تے \Log{N1i} لئی تھاں بدلّن تے تقلیل والیاں تبدیلیاں سدھیاں \Log{N2i} وچ ورتیاں جا سکدیاں نیں۔ ایس لئی نارمل روپ وچ لیاؤن دا قضیہ \Log{N2i} لئی وی سچ اے۔

\begin{thm}\ollabel{thm:normalization-N2i}
کوئی وی \Log{N2i}-!!{proof} $\delta$ نارمل روپ وچ آ سکدا اے، یعنی قطعاں توں بناں !!a{proof}~$\delta'$ وچ بدلیا جا سکدا اے۔
\end{thm}

\end{document}
